\documentclass{amsart}
% For dealing with references we use the comment environment
\usepackage{verbatim}
\newenvironment{reference}{\comment}{\endcomment}
%\newenvironment{reference}{}{}
\newenvironment{slogan}{\comment}{\endcomment}
\newenvironment{history}{\comment}{\endcomment}

% For commutative diagrams we use Xy-pic
\usepackage[all]{xy}

% We use 2cell for 2-commutative diagrams.
\xyoption{2cell}
\UseAllTwocells

% We use multicol for the list of chapters between chapters
\usepackage{multicol}

% This is generally recommended for better output
\usepackage{lmodern}
\usepackage[T1]{fontenc}

% For cross-file-references

% Package for hypertext links:
\usepackage{hyperref}

% For any local file, say "hello.tex" you want to link to please

% Theorem environments.
%
\theoremstyle{plain}
\newtheorem{theorem}[subsection]{Theorem}
\newtheorem{proposition}[subsection]{Proposition}
\newtheorem{lemma}[subsection]{Lemma}

\theoremstyle{definition}
\newtheorem{definition}[subsection]{Definition}
\newtheorem{example}[subsection]{Example}
\newtheorem{exercise}[subsection]{Exercise}
\newtheorem{situation}[subsection]{Situation}

\theoremstyle{remark}
\newtheorem{remark}[subsection]{Remark}
\newtheorem{remarks}[subsection]{Remarks}

\numberwithin{equation}{subsection}

% Macros
%
\def\lim{\mathop{\mathrm{lim}}\nolimits}
\def\colim{\mathop{\mathrm{colim}}\nolimits}
\def\Spec{\mathop{\mathrm{Spec}}}
\def\Hom{\mathop{\mathrm{Hom}}\nolimits}
\def\Ext{\mathop{\mathrm{Ext}}\nolimits}
\def\SheafHom{\mathop{\mathcal{H}\!\mathit{om}}\nolimits}
\def\SheafExt{\mathop{\mathcal{E}\!\mathit{xt}}\nolimits}
\def\Sch{\mathit{Sch}}
\def\Mor{\mathop{\mathrm{Mor}}\nolimits}
\def\Ob{\mathop{\mathrm{Ob}}\nolimits}
\def\Sh{\mathop{\mathit{Sh}}\nolimits}
\def\NL{\mathop{N\!L}\nolimits}
\def\CH{\mathop{\mathrm{CH}}\nolimits}
\def\proetale{{pro\text{-}\acute{e}tale}}
\def\etale{{\acute{e}tale}}
\def\QCoh{\mathit{QCoh}}
\def\Ker{\mathop{\mathrm{Ker}}}
\def\Im{\mathop{\mathrm{Im}}}
\def\Coker{\mathop{\mathrm{Coker}}}
\def\Coim{\mathop{\mathrm{Coim}}}

% Boxtimes
%
\DeclareMathSymbol{\boxtimes}{\mathbin}{AMSa}{"02}

%
% Macros for moduli stacks/spaces
%
\def\QCohstack{\mathcal{QC}\!\mathit{oh}}
\def\Cohstack{\mathcal{C}\!\mathit{oh}}
\def\Spacesstack{\mathcal{S}\!\mathit{paces}}
\def\Quotfunctor{\mathrm{Quot}}
\def\Hilbfunctor{\mathrm{Hilb}}
\def\Curvesstack{\mathcal{C}\!\mathit{urves}}
\def\Polarizedstack{\mathcal{P}\!\mathit{olarized}}
\def\Complexesstack{\mathcal{C}\!\mathit{omplexes}}
% \Pic is the operator that assigns to X its picard group, usage \Pic(X)
% \Picardstack_{X/B} denotes the Picard stack of X over B
% \Picardfunctor_{X/B} denotes the Picard functor of X over B
\def\Pic{\mathop{\mathrm{Pic}}\nolimits}
\def\Picardstack{\mathcal{P}\!\mathit{ic}}
\def\Picardfunctor{\mathrm{Pic}}
\def\Deformationcategory{\mathcal{D}\!\mathit{ef}}

\setlength{\emergencystretch}{3em}
\begin{document}
\title[Stacks Project, Tag 0FXI]{356 --- Strong generation by the ring forces finite-dimensional regularity}
\maketitle
\noindent Copyright (C) 2005--2025 Johan de Jong. Stacks Project authors. Source: \href{https://stacks.math.columbia.edu/tag/0FXI}{Tag 0FXI}.

\noindent Editorial difficulty note (not author text). Difficulty H2: The proof realizes a universal Ext class by a truncated finite free resolution and factors it through a sequence of cohomologically zero maps. The ghost lemma annihilates the composite, forcing a uniform projective-dimension bound for all finite modules and hence finite-dimensional regularity..

\noindent Editorial provenance note (not author text). History: excerpt from revision \texttt{a0444\allowbreak 6e57e\allowbreak c1fbc\allowbreak 252a8\allowbreak 71afc\allowbreak ec775\allowbreak 2fb28\allowbreak 07b14}, more-algebra.tex, lines 20710--20764. Reference presentation was adapted; original-author context and core support blocks were retained or added. The mathematical wording is unchanged. Licensed under GNU FDL 1.2 or later; no invariant sections or cover texts. See ../licenses/COPYING. Context blocks reproduce author definitions and supporting results; they are not additional counted problems.

\begin{lemma}
\label{lemma-not-regular-not-strong}
Let $R$ be a Noetherian ring. If $R$ is
a strong generator for $D_{perf}(R)$, then $R$ is regular
of finite dimension.
\end{lemma}

\begin{proof}
Assume $D_{perf}(R) = \langle R \rangle_n$ for some $n \geq 1$.
For any finite $R$-module $M$ we can choose a complex
$$
P = (
P^{-n - 1} \xrightarrow{d^{-n - 1}}
P^{-n} \xrightarrow{d^{-n}}
P^{-n + 1} \xrightarrow{d^1}
\ldots \xrightarrow{d^{-1}} P^0)
$$
of finite free $R$-modules with $H^i(P) = 0$ for $i = -n, \ldots, - 1$
and $M \cong \Coker(d^{-1})$. Note that $P$ is in $D_{perf}(R)$.
For any $R$-module $N$ we can compute $\Ext^n_R(M, N)$
the finite free resolution $P$ of $M$, see
Algebra, Section \href{https://stacks.math.columbia.edu/tag/00LO}{section ext (Tag 00LO)} and compare with
Derived Categories, Section \href{https://stacks.math.columbia.edu/tag/06XP}{section ext (Tag 06XP)}.
In particular, the sequence above defines an element
$$
\xi \in \Ext^n_R(\Coker(d^{-1}), \Coker(d^{-n - 1})) =
\Ext^n_R(M, \Coker(d^{-n - 1}))
$$
and for any element $\overline{\xi}$ in $\Ext^n_R(M, N)$
there is a $R$-module map $\varphi : \Coker(d^{-n - 1}) \to N$ such
that $\varphi$ maps $\xi$ to $\overline{\xi}$.
For $j = 1, \ldots, n - 1$ consider the complexes
$$
K_j = (\Coker(d^{-n - 1}) \to P^{-n + 1} \to \ldots \to P^{-j})
$$
with $\Coker(d^{-n - 1})$ in degree $-n$ and $P^t$ in degree $t$.
We also set $K_n = \Coker(d^{-n - 1})[n]$. Then we have maps
$$
P \to K_1 \to K_2 \to \ldots \to K_n
$$
which induce vanishing maps on cohomology. By Lemma \href{https://stacks.math.columbia.edu/tag/0FXH}{lemma ghost lemma (Tag 0FXH)}
since $P \in D_{perf}(R) = \langle R \rangle_n$ we find that the composition
of this maps is zero in $D(R)$. Since
$\Hom_{D(R)}(P, K_n) = \Hom_{K(R)}(P, K_n)$ by
Derived Categories, Lemma \href{https://stacks.math.columbia.edu/tag/064B}{lemma morphisms from projective complex (Tag 064B)}
we conclude $\xi = 0$. Hence $\Ext^n_R(M, N) = 0$ for all $R$-modules
$N$, see discussion above.
It follows that $M$ has projective dimension $\leq n - 1$ by
Algebra, Lemma \href{https://stacks.math.columbia.edu/tag/065R}{lemma projective dimension ext (Tag 065R)}.
Since this holds for all finite $R$-modules $M$ we conclude that
$R$ has finite global dimension, see
Algebra, Lemma \href{https://stacks.math.columbia.edu/tag/065T}{lemma finite gl dim (Tag 065T)}.
We finally conclude by Algebra, Lemma
\href{https://stacks.math.columbia.edu/tag/00OE}{lemma finite gl dim finite dim regular (Tag 00OE)}.
\end{proof}
\end{document}
